% !TEX root = ../../hw04.tex
% Author's handwritten solution: images/HW05/IMG_9282--IMG_9283.HEIC.
% Preserve the root-test and squeeze methods. Approved additions cover negative
% integer k, distinguish the root-test value from the radius, correct the
% reversed p-series conditions, and state the absolute-convergence set.
\begin{proof}[Solution]
  Write the power series as $\sum_{n=1}^{\infty}c_nx^n$, where
  \[
    c_n=\frac{1}{3^{n/2}n^k}.
  \]
  We use the root test. For every $n\ge1$,
  \[
    \left|\frac{x^n}{(\sqrt3)^n n^k}\right|^{1/n}
    =\frac{|x|}{\sqrt3\,(n^k)^{1/n}}
    =\frac{|x|}{\sqrt3\,(n^{1/n})^k}.
  \]
  First suppose $k\ge0$.
  Choose a positive integer $m\ge k$. Since
  $n^{1/n}\to1$, we have
  \[
    1\le(n^{1/n})^k\le(n^{1/n})^m\longrightarrow1.
  \]
  The squeeze theorem gives $(n^{1/n})^k\to1$.

  If $k<0$, apply the same argument to the positive integer $-k$.
  Taking reciprocals gives
  \[
    (n^{1/n})^k=\frac{1}{(n^{1/n})^{-k}}\longrightarrow1.
  \]
  Thus this limit holds for every integer $k$.

  It follows that
  \[
    \lim_{n\to\infty}
      \left|\frac{x^n}{(\sqrt3)^n n^k}\right|^{1/n}
    =\frac{|x|}{\sqrt3}.
  \]
  We obtain absolute convergence when
  \[
    \frac{|x|}{\sqrt3}<1
    \quad\Longleftrightarrow\quad
    |x|<\sqrt3
    \quad\Longleftrightarrow\quad
    -\sqrt3<x<\sqrt3.
  \]
  If $|x|>\sqrt3$, the series diverges
  because its terms do not tend to zero, by the root
  divergence argument.
  Hence the radius of convergence is $\sqrt3$;
  the value $|x|/\sqrt3$ above is the root-test limit, not the radius.

  At the endpoints $x=\pm\sqrt3$, the absolute-value series becomes
  \[
    \sum_{n=1}^{\infty}
      \left|\frac{(\pm\sqrt3)^n}{(\sqrt3)^n n^k}\right|
    =\sum_{n=1}^{\infty}\frac{1}{n^k}.
  \]
  This series converges if $k>1$ and diverges if
  $k\le1$.
  For $k=1$ it is the harmonic series. For $k\le0$
  its terms do not tend to zero. For $k>1$ the p-series test applies.
  Thus both endpoints give absolute convergence exactly when $k>1$.

  The set of absolute convergence is therefore
  \[
    \begin{cases}
      [-\sqrt3,\sqrt3],&k>1,\\
      (-\sqrt3,\sqrt3),&k\le1.
    \end{cases}
  \]

  Failure of absolute convergence does not always imply divergence
  of the original series.
  For example, at $x=-\sqrt3$ and $k=1$ the
  alternating harmonic series converges conditionally.
\end{proof}
